\begin{figure}[htbp]
\centering
\begin{tikzpicture}[scale=1.3]
 \fill[blue!9] (-.675,-1.5)--(-.21,1.6)--(1.65,1.6)--cycle;
 \draw[blue!65!black,thick] (-.675,-1.5)--(-.21,1.6)--(1.65,1.6)--cycle;
 \fill[orange!30] (-.315,-.3)--(-.09,0)--(.15,1.6)--(1.29,1.6)--(.09,0)--(.045,-.3)--cycle;
 \draw[orange!70!black] (-.315,-.3)--(-.09,0)--(.15,1.6)--(1.29,1.6)--(.09,0)--(.045,-.3)--cycle;
 \draw[dashed] (-.675,-1.5)--(.72,1.6);
 \draw[-{Latex}] (-1,0)--(1.95,0) node[right] {$v$};
 \draw[-{Latex}] (0,-1.7)--(0,1.93) node[above] {$u$};
 \node[left] at (-.68,-1.45) {$-Cd_I$};
 \node[left] at (-.34,-.3) {$-Ad_I$};
 \node[right] at (1.65,1.6) {$R_I+Ad_I$};
 \node[blue!65!black] at (1.1,.45) {$T_I$};
\end{tikzpicture}
\caption{In intercept coordinates, the orange polygon bounds the
frozen sweep, including the small region below $u=0$. The enclosing
blue triangle extends its tip to $u=-Cd_I$; the dashed line is
$v=m_Iu$. The drawing is schematic; the triangle and footprint are
defined by \eqref{eq:triangle} and \eqref{eq:old-footprint}.}
\label{fig:triangle}
\end{figure}
